\chapter{Introduction}
\label{ch:introduction}

% ---------------------------------------------------------------------------
\section{Motivation and Problem}
\label{sec:motivation}

Standing up production-grade cloud infrastructure requires expertise that
most engineering teams don't have all in one place. A backend developer
who can design a distributed system from scratch may have never written an
Nginx upstream block; a DevOps engineer fluent in Terraform may not know
how to construct a valid Helm chart for a stateful Kubernetes workload.
The result is that deploying even a moderately complex application to AWS
requires simultaneous competence across Terraform, Kubernetes, Helm, Nginx,
Jenkins, Ansible, and Prometheus, each with its own syntax, security model,
and failure modes. Small teams and time-constrained projects routinely
either skip parts of the stack entirely or write configurations at the edge
of someone's knowledge, under deadline pressure.

The more consequential problem is security. Infrastructure-as-Code files
are written once and trusted indefinitely. A Terraform module that grants
overly broad IAM permissions on the first day of a project will still be
granting them a year later, long after the engineer who wrote it has moved
on. Tools like Checkov and tfsec exist precisely because this pattern is
widespread enough to warrant dedicated static analysis tooling with hundreds
of distinct checks, each corresponding to a class of vulnerability that human
authors regularly introduce. But these tools are reactive: they report what
is wrong; they don't fix it, and they don't prevent the next engineer from
making the same mistake.

Existing automation closes part of the gap. Terraform modules, Helm charts,
and cloud provider templates reduce repetition, but they require the engineer
to already know which components to assemble, how to compose them correctly,
and how to harden the result. They are reusable parts, not a reasoning
system. An engineer who doesn't know what they need can't ask a module
registry for it in plain English; a module that deploys correctly isn't
necessarily one that deploys securely.

The problem this thesis addresses is not that IaC tooling is insufficient.
The problem is that the authorship and security validation steps sit outside
the reach of engineers who lack breadth across the full stack, and no
existing tool bridges that gap end to end.

% ---------------------------------------------------------------------------
\section{Goal, Scope, and Success Criteria}
\label{sec:goal}

The goal of this thesis is to design and implement InfraArchAgent: a
multi-agent system that accepts a natural-language description of what an
engineer wants to deploy on AWS and returns a security-scanned candidate
infrastructure package, without requiring the engineer to write a single
line of IaC.

\subsection*{In scope}

\begin{itemize}
  \item A four-stage pipeline with six agent instances: ArchitectAgent
    (plan generation), three parallel GeneratorAgents (IaC file generation
    optimised for cost, performance, and security), SecurityAgent (iterative
    scanning and autonomous remediation using Checkov and tfsec), and
    ValidatorAgent (cross-file consistency checking).
  \item A React frontend with real-time pipeline progress via Server-Sent
    Events, side-by-side variant comparison, security diff view, and ZIP
    download.
  \item A FastAPI backend with a PostgreSQL persistence layer storing
    pipeline runs, generated packages, and agent events.
  \item A provider-agnostic LLM interface supporting both the Anthropic
    Messages API and the OpenAI chat completions API, implemented via
    the Adapter pattern.
  \item AWS as the sole target cloud provider.
\end{itemize}

\subsection*{Non-goals}

The system does not execute or apply generated infrastructure against any
cloud provider. It holds no cloud credentials. It does not support Azure
or GCP in the current implementation, though the architecture permits
extension. It does not provide cost estimation, live infrastructure
monitoring, or state management. It is a single-user prototype; multi-user
authentication and run ownership are out of scope.

\subsection*{Success criteria}

\begin{itemize}
  \item A full pipeline run completes within 5 minutes from input submission
    to all available validated packages appearing in the UI, measured as
    the median of 10 runs against a real LLM API (PR-01).
  \item The SecurityAgent reduces the violation count on every generated
    package; packages with unresolved violations after the configured
    iteration limit are delivered with violations clearly flagged, giving
    the engineer full visibility into what requires manual attention.
  \item The UI reflects each agent state change within 1 second of the
    backend emitting the corresponding SSE event (PR-02).
  \item Adding a new LLM provider requires implementing exactly one adapter
    class with two methods and no changes to any agent (NFR-03), verified
    by the stub-provider integration test.
\end{itemize}

% ---------------------------------------------------------------------------
\section{Three Principal Contributions}
\label{sec:contributions}

\begin{enumerate}

  \item \textbf{C1 --- Cooperating and competing multi-agent pipeline.}
    Coordinating agents that must run in a fixed sequence while embedding
    a parallel competition stage inside that sequence is an orchestration
    problem with no off-the-shelf solution. The three GeneratorAgents must
    run concurrently and independently, yet the overall pipeline must still
    respect the sequential dependency between plan generation and security
    scanning. This thesis designs and implements a Python asyncio
    orchestration layer that runs the three generators as concurrent tasks
    behind a barrier, handles partial failures without halting the pipeline,
    and keeps inter-agent communication decoupled through a shared
    IaCPackage data object. The provider-agnostic LLM interface, implemented
    via the Adapter pattern, supports this pipeline by making every agent
    independent of any specific model provider.
    Evidence: the sequence diagram in Chapter~\ref{ch:design}, the asyncio
    implementation in Chapter~\ref{ch:implementation}, and the concurrency
    tests in Chapter~\ref{ch:testing}.

  \item \textbf{C2 --- Autonomous security remediation loop.}
    Generating syntactically correct IaC is not the hard part; generating
    IaC that passes Checkov and tfsec without manual intervention is. The
    SecurityAgent must interpret structured scanner output, produce a
    targeted fix for each violation via an LLM call, apply that fix to the
    relevant file, and rescan, repeating until either all checks pass or the
    configured iteration limit is reached. Fixes can introduce new
    violations, so the loop must detect convergence, not just completion.
    This thesis implements the scan-remediate-rescan cycle with configurable
    iteration limits, per-iteration diff recording, and clean handling of
    packages that exhaust the limit without fully resolving.
    Evidence: FR-S-01 through FR-S-06 in Chapter~\ref{ch:requirements},
    the SecurityAgent design in Chapter~\ref{ch:design}, and the
    remediation integration tests in Chapter~\ref{ch:testing}.

  \item \textbf{C3 --- Interactive UI with real-time pipeline visualization.}
    Giving engineers visibility into a multi-agent pipeline as it runs
    requires streaming state updates from multiple concurrent async tasks
    to a browser without polling. Three agents run in parallel and each
    can emit state changes at any time; the UI must reflect those changes
    within one second without page refresh. This thesis implements a
    Server-Sent Events stream that serialises concurrent agent events
    through an asyncio queue and pushes them to a React frontend, with
    four purpose-built screens covering input, live progress, variant
    comparison, and security diff review. The comparison panel presents
    three generated variants side by side with file trees, security scores,
    and a unified diff of every change the SecurityAgent made, giving
    the engineer a complete picture before committing to a download.
    Evidence: FR-UI-01 through FR-UI-06 and FR-P-04 in
    Chapter~\ref{ch:requirements}, the UI design in
    Chapter~\ref{ch:design}, and the SSE integration tests in
    Chapter~\ref{ch:testing}.

\end{enumerate}

\EvidenceTable

% ---------------------------------------------------------------------------
\section{Thesis Roadmap}
\label{sec:roadmap}

Chapter~\ref{ch:requirements} establishes what InfraArchAgent must do,
formalising the pipeline behaviour, interface contracts, performance targets,
and use cases as testable requirements. The functional requirements table
and the traceability matrix connect every commitment in this chapter to
design evidence, code, and tests in later chapters.

Chapter~\ref{ch:design} explains the structural decisions that satisfy
those requirements: the four-layer architecture, the class hierarchy that
encodes the Strategy and Adapter patterns, the PostgreSQL schema with its
JSONB package storage, the SSE-based real-time update mechanism, and the
six architectural decisions with their rejected alternatives.

Chapter~\ref{ch:implementation} documents the actual construction: the
asyncio orchestration model covering concurrent programming, the
object-oriented class hierarchy covering OOM and OOD, the PostgreSQL schema
and query design covering the database examination topic, and the integration
of Checkov and tfsec as subprocesses.

Chapter~\ref{ch:testing} provides evidence that the system works: unit
tests for each agent in isolation, integration tests for the full pipeline,
performance benchmarks against the 5-minute completion target, and security
remediation tests against known-vulnerable Terraform templates.

Chapter~\ref{ch:conclusion} reflects on what was built, what the
measurements showed, where the current implementation falls short, and
what a follow-on project would address first.